\documentclass[10pt,a4paper]{amsart}
\usepackage[margin=19mm]{geometry}
\usepackage{amsmath,amssymb,mathtools}
\usepackage[scr=rsfs,cal=euler]{mathalfa}
\usepackage{xfrac}
\usepackage[unicode,hidelinks,bookmarks=false]{hyperref}
\newcommand{\ie}{i.e.\ }
\newcommand{\cf}{cf.\ }
\newcommand{\resp}{respectively\ }
\newcommand{\Subsection}{\subsection}
\providecommand{\nobreakdash}{\nobreak\hbox{-}}
\setlength{\emergencystretch}{2em}
\multlinegap0pt
\usepackage{bm}
\usepackage{bbm}
\usepackage{dsfont}
\usepackage{xspace}
\usepackage{cancel}
\usepackage{mathtools}
\usepackage[shortlabels]{enumitem}
\usepackage{amsthm}
\makeatletter
\@ifundefined{theorem}{\newtheorem{theorem}{Theorem}}{}
\@ifundefined{claim}{\newtheorem{claim}[theorem]{Claim}}{}
\@ifundefined{lemma}{\newtheorem{lemma}[theorem]{Lemma}}{}
\@ifundefined{observation}{\newtheorem{observation}[theorem]{Observation}}{}
\makeatother
\title[Hindrance from a wasteful partial linkage]{Hindrance from a wasteful partial linkage}
\author{Attila Jo\'o}
\date{Source: arXiv:2410.19583v2 (CC BY 4.0)}
\begin{document}
\maketitle

\noindent\emph{Source and licence.} Hindrance from a wasteful partial linkage. Version arXiv:2410.19583v2 (CC BY 4.0), distributed under the Creative Commons Attribution 4.0 International licence, as recorded in the arXiv licence metadata for this exact version and shown on the abstract page https://arxiv.org/abs/2410.19583v2. Full rights notice: licenses/arxiv-2410.19583-CC-BY-4.0.txt.\par\smallskip

\noindent\textit{Editorial scope.} Counted unit: the Lemma whose source label is \texttt{lem: key lemma} in the article (source0.tex lines 500--507, source label \texttt{lem: key lemma}), delivered together with the article's own complete original proof of it (source0.tex lines 508--570, one \texttt{proof} environment of 63 lines). No mathematics is written, repaired or re-derived: statement and proof are the article's own text line for line. Required context retained uncounted: obs: dis to strongly dis (lines 254--258); lem: remains alternating trail (lines 450--459). Every definition, display and label of those lines is retained unchanged. Omitted from this package and not redistributed: 1--253, 259--449, 460--499, 571--802. Nothing inside a retained excerpt is omitted or altered: every retained body is delivered exactly as the article's lines are written, so no omission receipt of any kind accompanies this package. Citations: the retained excerpts carry no citation command at all, so no bibliography entry of the article is transcribed and no reference is redistributed. Reference adaptation: none was needed and none was made; every reference inside the delivered unit resolves inside it, so no unresolved reference or citation is produced. Printed slips kept literal: no printed word, symbol, hypothesis or number of the counted statement or of its proof was altered. Layout and class: the article's own class is \texttt{amsart} with packages including amsmath, amssymb, amsthm, bm, mathtools, tikz, thmtools, float, xcolor, hyperref, biblatex, inputenc, csquotes, fontenc; the wrapper is the lane's pinned \texttt{amsart} editorial wrapper at 10pt on A4, which restates the theorem-like environments the retained text needs and adds the article's own macro definitions that the retained text needs (none), with extra wrapper packages (bm, bbm, dsfont, xspace, cancel, mathtools, enumitem). The delivered numbering is the unit's own. Assets: no retained span contains a figure environment, a graphics-inclusion command, a PDF-page inclusion, a table or a TikZ picture; no image file is copied into this package, the package declares no asset dependency and no image file is redistributed with it. Licence: version arXiv:2410.19583v2 of this article is distributed under the Creative Commons Attribution 4.0 International licence, as recorded in the arXiv licence metadata for this exact version and shown on the abstract page https://arxiv.org/abs/2410.19583v2. A full rights notice is delivered at licenses/arxiv-2410.19583-CC-BY-4.0.txt. Characters: every retained excerpt is pure ASCII, so the delivered engine, XeLaTeX with its default Latin Modern text font, reports no missing character.
\begin{observation}\label{obs: dis to strongly dis}
For every infinite cardinal $ \kappa $, every $ \kappa $-set of disjoint alternating trails includes a $ \kappa $-subset of 
strongly 
disjoint alternating trails.
\end{observation}

\begin{lemma}\label{lem: remains alternating trail}
For every $ n<\omega $, the following holds:
\begin{enumerate}
\item Every $ \mathcal{L}_{n} $-alternating trail $ T $ with 
$ V(T)\cap (W_n \cup N_{D_{n}}^{-}(U_n))\subseteq \{ \mathsf{ter}(T) \} $ is an $ \mathcal{L}_{n+1} $-alternating 
trail.
\item Every $ \mathcal{L}_{n+1} $-alternating trail $ T $ with $ V(T)\cap W_n \subseteq \{ \mathsf{ter}(T) \} $ is an $ 
\mathcal{L}_n $-alternating trail.
\end{enumerate}
\end{lemma}

\begin{lemma}\label{lem: key lemma}
For every $ n<\omega $:
\begin{enumerate}[label=(\roman*)]
\item[$ (i_n) $]\label{item: no kappa aug} there is no $ \kappa $-set $ \mathcal{T} $ of disjoint $ \mathcal{L}_n $-augmenting trails;
\item[$ (ii_n) $]\label{item: few meets Wn} there is no $ \kappa $-set $ \mathcal{T} $ of disjoint $ \mathcal{L}_n 
$-alternating $ \widehat{A}W_n$-trails.
\end{enumerate}
\end{lemma}

\begin{proof}
We apply induction on $ n $. Clearly, $ (i_0) $ holds because $ \left|\widehat{B}_0\right|<\kappa $ (by $ 
\left|\widehat{B}\right|<\left|\widehat{A}\right|=\kappa $) and every $ 
\mathcal{L}_0 $-augmenting trail
terminates in $ 
 \widehat{B}_0 $ by definition. Suppose that we already know $ (i_n) $ for some $ n $.   Next we prove $ (ii_n) $ and $ 
  (i_{n+1})$.
\begin{claim}\label{clm: no lot trail to Un}
 There is no $ \kappa $-set of $ \mathcal{L}_n 
  $-augmenting trails with $ \mathsf{ter}(\mathcal{T})\subseteq U_n $ where any two trails in $ \mathcal{T} $ are either disjoint or sharing 
   exactly their terminal vertex.
\end{claim}
\begin{proof}
If $ \left|\mathsf{ter}(\mathcal{T}) \right|<\kappa$, then, by the regularity of $ \kappa $, there is 
 a 
 $ v\in \mathsf{ter}(\mathcal{T}) $ such that $ \kappa $ many $ T\in \mathcal{T} $ terminates at $ v $, contradicting  $v\in U_n $. Thus $ 
 \left|\mathsf{ter}(\mathcal{T})\right|=\kappa $. But then by choosing for each $ v\in \mathsf{ter}(\mathcal{T}) $ exactly one $ T\in 
 \mathcal{T} $ with $ \mathsf{ter}(T)=v $, we obtain a $ \kappa $-set of disjoint $ \mathcal{L}_n 
 $-augmenting trails which contradicts $ (i_n) $.
\end{proof}
 
 
 

 
 Suppose for a contradiction that $ (ii_n) $ fails, i.e. there is a $ \kappa $-set $ \mathcal{T} $ of disjoint $ \mathcal{L}_n 
 $-alternating $ 
 \widehat{A}W_n $-trails. By Observation \ref{obs: dis to strongly dis}, we may assume that the trails in $ \mathcal{T} $ 
 are strongly disjoint. For $ 
 T\in \mathcal{T} $, let 
 $ P_T $ be the unique path in $ \mathcal{P}_n $ with $ \mathsf{ter}(T)\in V(P_T) $ and let $ v_P $ be the terminal vertex of 
 the unique path in $ \mathcal{P}_{n+1} $ that is a proper initial segment of $ P_T $ (exists by the definition of $ W_n $). 
 Let $ T' $ be a forward-extension of $ 
 T $ that 
 we obtain by going backwards along $ P_T $ until  $ v_P $  and then adding an outgoing edge of $ v_P $ in $ D_n $
 with head in $ U_n $ 
 (which exists because $ v_P $ is nothing else than the first vertex of $ P_T $ that is in $ N_{D_{n}}^{-}(U_n) $). Then $ 
 \mathcal{T}':=\{ T':\ T\in \mathcal{T} \} $ is a $ \kappa $-set of $ \mathcal{L}_n 
 $-augmenting paths with $ \mathsf{ter}(\mathcal{T}')\subseteq U_n $ where any two trails in $ \mathcal{T}' $ are either disjoint or sharing 
 exactly their terminal vertex which contradicts Claim \ref{clm: no lot trail to Un}.
 
 Suppose for a contradiction that $ (i_{n+1}) $ fails, i.e. there is a $ \kappa $-set $ \mathcal{T} $ of disjoint $ 
 \mathcal{L}_{n+1} 
 $-augmenting trails. By $( i_n) $ we know that less than $ \kappa $ many $ T\in \mathcal{T} $ are $ \mathcal{L}_n 
 $-augmenting trails, thus by 
 deleting these we can assume that no $ T\in \mathcal{T} $ is an $ \mathcal{L}_n $-augmenting trail. then exactly one of the following holds for 
 each
  $ T\in \mathcal{T} $: either $ T $ is not even an $ \mathcal{L}_n $-alternating trail or it is but $ 
 \mathsf{ter}(T)\in B_{n+1}\setminus B_n $. Since one of these two cases occurs $ \kappa $ often, we can assume that one of them holds 
 actually for every $ T\in \mathcal{T} $. Suppose first, that no trail in $ \mathcal{T} $ is $ \mathcal{L}_n $-alternating. 
 Then  (2) of Lemma 
 \ref{lem: remains alternating trail}   guarantees that each $ 
 T\in \mathcal{T} $ meets $ W_n $. Let $ \mathcal{T}' $ consists of the initial segments of the trails in $ \mathcal{T} $ up to their first vertex in $ 
 W_n $. Then (2) 
 of Lemma \ref{lem: remains alternating trail} ensures that $ 
 \mathcal{T}' 
 $ consists of $ \mathcal{L}_n $-alternating trails. Since $ \left|\mathcal{T}'\right|=\kappa $, this contradicts $ (ii_n) $. Assume now that every 
 $ T\in \mathcal{T} $ is $ \mathcal{L}_n $-alternating but terminates in $ B_{n+1}\setminus B_n= N_{D_n}^{-}(U_n) $. Then 
 every $ T\in 
 \mathcal{T} $ can be extended forward by a new edge to  an $ \mathcal{L}_n $-augmenting trail $ T' $ that terminates in $ 
 U_n $. But then the set 
 $ \mathcal{T}' $  of these extensions contradicts Claim \ref{clm: no lot trail to Un}.
\end{proof}

\end{document}
